% ============================================================
\chapter{Deep Q-Networks and Variants}
\label{ch:dqn}
% ============================================================

In February 2015, a team from DeepMind published in \emph{Nature} a result that made headlines: an algorithm, the \emph{Deep Q-Network} (DQN), learned to play 49 Atari games directly from screen pixels, achieving superhuman performance in several of them. The key: combining the tabular Q-learning of the previous chapter with a deep neural network that approximates the $Q$ function. Two technical innovations---\emph{experience replay} (storing transitions and resampling them randomly) and the \emph{target network} (updating targets more slowly)---stabilised a learning process that, without them, diverges. DQN opened the era of \emph{deep reinforcement learning}.

\begin{intuition}
The Deep Q-Network (DQN) by Mnih et al.\ (2013, 2015) revolutionized RL
by combining Q-Learning with deep neural networks. Two key innovations
--- experience replay and a target network --- stabilize training,
enabling agents to learn directly from raw pixels and achieve human-level
performance on Atari games.
\end{intuition}

% ------------------------------------------------------------
\section{From Tabular Q to Function Approximation}
\label{sec:dqn:approx}
% ------------------------------------------------------------

Tabular Q-Learning cannot handle high-dimensional or continuous state
spaces. The idea is to approximate $Q^*(s,a)$ with a neural network
parameterized by $\theta$:
\[
    Q(s, a; \theta) \approx Q^*(s, a).
\]

\begin{attention}[title=\textbf{Naive Instability}]
\index{deadly triad}
Directly applying Q-Learning with a neural network diverges due to three
issues (the \textbf{deadly triad}):
\begin{enumerate}
    \item \textbf{Temporal correlation}: consecutive samples
    $(s_t, a_t, r_{t+1}, s_{t+1})$ are strongly correlated.
    \item \textbf{Non-stationary target}: the target
    $r + \gamma \max_{a'} Q(s', a'; \theta)$ depends on $\theta$,
    which changes at every update.
    \item \textbf{Function approximation}: Q-Learning convergence
    is no longer guaranteed with nonlinear approximators.
\end{enumerate}
\end{attention}

% ------------------------------------------------------------
\section{DQN Architecture}
\label{sec:dqn:arch}
% ------------------------------------------------------------

\begin{definition}[Deep Q-Network]
\index{DQN}
The \textbf{DQN} uses two stabilization mechanisms:
\begin{enumerate}
    \item \textbf{Experience replay}\index{experience replay}: transitions
    $(s, a, r, s')$ are stored in a replay buffer $\mathcal{D}$ of capacity
    $N$ and uniformly sampled in mini-batches for updates.
    \item \textbf{Target network}\index{target network}: a separate network
    $Q(s,a;\theta^-)$ with frozen parameters $\theta^-$ computes the target.
    Parameters $\theta^-$ are copied from $\theta$ every $C$ steps.
\end{enumerate}
\end{definition}

\begin{keyformulas}[title=\textbf{DQN Loss Function}]
The loss is the mean squared error over a mini-batch sampled from
$\mathcal{D}$:
\[
    \mathcal{L}(\theta) = \E_{(s,a,r,s') \sim \mathcal{D}} \left[
    \left( r + \gamma \max_{a'} Q(s', a'; \theta^-) - Q(s, a; \theta) \right)^2
    \right].
\]
\end{keyformulas}

\begin{algorithme}[title=\textbf{DQN Algorithm}]
\begin{enumerate}
    \item Initialize replay buffer $\mathcal{D}$ with capacity $N$
    \item Initialize $Q(s,a;\theta)$ with random weights $\theta$
    \item Set target parameters $\theta^- \leftarrow \theta$
    \item For each episode:
    \begin{enumerate}
        \item Initialize state $s_0$ (preprocess frame stack)
        \item For each step $t$:
        \begin{enumerate}
            \item Select $a_t = \argmax_a Q(s_t, a; \theta)$ with probability
                  $1-\epsilon$, else random action
            \item Execute $a_t$, observe $r_{t+1}, s_{t+1}$
            \item Store $(s_t, a_t, r_{t+1}, s_{t+1})$ in $\mathcal{D}$
            \item Sample mini-batch from $\mathcal{D}$
            \item Compute targets: $y_j = r_j + \gamma \max_{a'} Q(s'_j, a'; \theta^-)$
            \item Update $\theta$ by gradient descent on
                  $\frac{1}{B}\sum_j (y_j - Q(s_j, a_j; \theta))^2$
            \item Every $C$ steps: $\theta^- \leftarrow \theta$
        \end{enumerate}
    \end{enumerate}
\end{enumerate}
\end{algorithme}

\begin{theorem}[Experience Replay Decorrelation]
Uniform sampling from a sufficiently large replay buffer produces
approximately i.i.d.\ samples, satisfying the stochastic gradient
descent assumption. Each transition can be reused multiple times,
improving sample efficiency.
\end{theorem}

% ------------------------------------------------------------
\section{DQN for Atari}
\label{sec:dqn:atari}
% ------------------------------------------------------------

\begin{example}[Atari 2600 Results]
\index{Atari}
The DQN architecture for Atari processes a stack of 4 grayscale
$84 \times 84$ frames through three convolutional layers followed by
two fully connected layers. Key results from Mnih et al.\ (2015):
\begin{itemize}
    \item Superhuman performance on 29 out of 49 Atari games.
    \item Single architecture and hyperparameters across all games.
    \item Learned from raw pixels with reward clipping to $[-1, +1]$.
\end{itemize}
\end{example}

% ------------------------------------------------------------
\section{Double DQN}
\label{sec:ddqn}
% ------------------------------------------------------------

\begin{definition}[Double DQN]
\index{Double DQN}
Double DQN (van Hasselt et al., 2016) addresses the maximization bias
by decoupling action selection from evaluation:
\[
    y_t = R_{t+1} + \gamma Q\Big(S_{t+1},
    \argmax_{a'} Q(S_{t+1}, a'; \theta)\,; \theta^-\Big).
\]
The online network $\theta$ selects the best action; the target network
$\theta^-$ evaluates it.
\end{definition}

\begin{proposition}[Overestimation Reduction]
Standard DQN overestimates Q-values because
$\E[\max_a Q] \geq \max_a \E[Q]$. Double DQN reduces this bias
significantly, leading to more stable training and improved final
performance on most Atari games.
\end{proposition}

% ------------------------------------------------------------
\section{Dueling DQN}
\label{sec:dueling}
% ------------------------------------------------------------

\begin{definition}[Dueling Architecture]
\index{Dueling DQN}
The dueling architecture (Wang et al., 2016) decomposes the Q-function:
\[
    Q(s, a; \theta, \alpha, \beta) = V(s; \theta, \beta)
    + \left( A(s, a; \theta, \alpha) - \frac{1}{\abs{\mathcal{A}}}
    \sum_{a'} A(s, a'; \theta, \alpha) \right).
\]
\end{definition}

\begin{intuition}
The value stream $V(s)$ captures state quality independently of action
choice. The advantage stream $A(s,a)$ captures the \emph{relative}
benefit of each action. This decomposition accelerates learning in
states where action choice has little impact.
\end{intuition}

% ------------------------------------------------------------
\section{Prioritized Experience Replay}
\label{sec:per}
% ------------------------------------------------------------

\begin{definition}[Prioritized Replay]
\index{prioritized experience replay}
Prioritized experience replay (Schaul et al., 2016) samples transitions
with probability proportional to their TD error magnitude:
\[
    P(i) = \frac{p_i^\alpha}{\sum_k p_k^\alpha}, \qquad
    p_i = \abs{\delta_i} + \epsilon_p
\]
where $\delta_i$ is the last observed TD error, $\alpha \in [0,1]$
controls prioritization, and $\epsilon_p > 0$ ensures nonzero
probability.
\end{definition}

\begin{definition}[Importance Sampling Correction]
To compensate for the non-uniform sampling bias, importance sampling
weights are applied:
\[
    w_i = \left(\frac{1}{N \cdot P(i)}\right)^\beta
\]
where $\beta$ is annealed from $\beta_0$ to $1$ during training. The
weighted loss becomes
$\mathcal{L} = \frac{1}{B}\sum_j w_j (y_j - Q(s_j, a_j; \theta))^2$.
\end{definition}

\begin{remark}
Prioritized replay significantly improves sample efficiency by focusing
updates on ``surprising'' transitions. The sum-tree data structure
enables efficient $O(\log N)$ priority-based sampling.
\end{remark}

% ------------------------------------------------------------
\section{Rainbow DQN}
\label{sec:rainbow}
% ------------------------------------------------------------

\begin{definition}[Rainbow]
\index{Rainbow}
Rainbow (Hessel et al., 2018) combines six DQN extensions:
\begin{enumerate}
    \item Double Q-Learning
    \item Prioritized experience replay
    \item Dueling architecture
    \item Multi-step returns ($n$-step TD targets)
    \item Distributional RL (C51)
    \item Noisy networks (learned exploration)
\end{enumerate}
\end{definition}

\begin{keyformulas}[title=\textbf{$n$-Step DQN Target}]
\[
    y_t^{(n)} = \sum_{k=0}^{n-1} \gamma^k R_{t+k+1}
    + \gamma^n \max_{a'} Q(S_{t+n}, a'; \theta^-)
\]
\end{keyformulas}

\begin{proposition}[Rainbow Ablation]
Ablation studies show that prioritized replay and $n$-step returns
provide the largest individual contributions. Removing distributional
RL or dueling networks also hurts performance, but less dramatically.
\end{proposition}

% ------------------------------------------------------------
\section{Soft Target Updates}
\label{sec:soft_target}
% ------------------------------------------------------------

\begin{definition}[Polyak Averaging]
\index{Polyak averaging}
Instead of periodically copying $\theta \to \theta^-$, soft updates
blend the parameters at every step:
\[
    \theta^- \leftarrow \tau \theta + (1 - \tau) \theta^-
\]
with $\tau \ll 1$ (typically $\tau = 0.005$). This provides smoother
target evolution and is widely used in actor-critic methods.
\end{definition}

% ------------------------------------------------------------
\section{Python Implementation}
\label{sec:dqn:code}
% ------------------------------------------------------------

\begin{codeblock}[title=\textbf{Minimal DQN with PyTorch}]
\begin{minted}{python}
import torch
import torch.nn as nn
import numpy as np
from collections import deque
import random

class QNetwork(nn.Module):
    def __init__(self, state_dim, action_dim, hidden=128):
        super().__init__()
        self.net = nn.Sequential(
            nn.Linear(state_dim, hidden), nn.ReLU(),
            nn.Linear(hidden, hidden), nn.ReLU(),
            nn.Linear(hidden, action_dim)
        )

    def forward(self, x):
        return self.net(x)

class ReplayBuffer:
    def __init__(self, capacity=100000):
        self.buffer = deque(maxlen=capacity)

    def push(self, s, a, r, s_next, done):
        self.buffer.append((s, a, r, s_next, done))

    def sample(self, batch_size):
        batch = random.sample(self.buffer, batch_size)
        s, a, r, s2, d = zip(*batch)
        return (np.array(s), np.array(a), np.array(r),
                np.array(s2), np.array(d, dtype=np.float32))

def dqn_train(env, n_episodes=500, gamma=0.99, lr=1e-3,
              batch_size=64, target_update=100):
    sd = env.observation_space.shape[0]
    ad = env.action_space.n
    q_net = QNetwork(sd, ad)
    q_target = QNetwork(sd, ad)
    q_target.load_state_dict(q_net.state_dict())
    optimizer = torch.optim.Adam(q_net.parameters(), lr=lr)
    buf = ReplayBuffer()
    step = 0
    for ep in range(n_episodes):
        s, _ = env.reset()
        done = False
        while not done:
            eps = max(0.01, 1.0 - step / 5000)
            if random.random() < eps:
                a = env.action_space.sample()
            else:
                with torch.no_grad():
                    a = q_net(torch.FloatTensor(s)).argmax().item()
            s2, r, term, trunc, _ = env.step(a)
            buf.push(s, a, r, s2, term)
            s = s2
            done = term or trunc
            step += 1
            if len(buf.buffer) >= batch_size:
                sb, ab, rb, s2b, db = buf.sample(batch_size)
                q_vals = q_net(torch.FloatTensor(sb)).gather(
                    1, torch.LongTensor(ab).unsqueeze(1)).squeeze()
                with torch.no_grad():
                    q_next = q_target(
                        torch.FloatTensor(s2b)).max(1)[0]
                    targets = (torch.FloatTensor(rb)
                        + gamma * q_next
                        * (1 - torch.FloatTensor(db)))
                loss = nn.MSELoss()(q_vals, targets)
                optimizer.zero_grad()
                loss.backward()
                optimizer.step()
            if step % target_update == 0:
                q_target.load_state_dict(q_net.state_dict())
    return q_net
\end{minted}
\end{codeblock}

% ------------------------------------------------------------
\section{Exercises}
\label{sec:dqn:exercises}
% ------------------------------------------------------------

\begin{exercise}[DQN on CartPole]
Implement DQN for \texttt{CartPole-v1}. Train with and without a target
network. Plot the learning curves and explain the effect of the target
network on stability.
\end{exercise}

\begin{exercise}[Double DQN]
Modify the DQN implementation to use Double DQN. Compare Q-value
estimates with standard DQN on \texttt{LunarLander-v2} and verify that
Double DQN reduces overestimation.
\end{exercise}

\begin{exercise}[Dueling Architecture]
Implement the dueling architecture. Train on \texttt{CartPole-v1} and
compare with the standard architecture. In which states does the
advantage stream have the largest variance?
\end{exercise}

\begin{exercise}[Prioritized Replay]
Implement prioritized experience replay with importance sampling
correction. Compare sample efficiency with uniform replay on
\texttt{LunarLander-v2}.
\end{exercise}
